\chapter{Formas Normales Superiores (BCNF, 4FN, 5FN, DKNF) y Desnormalizaci\'on}
\label{cap:formas_superiores_desnormalizacion}

\section{Marco Conceptual y Te\'orico (Curr\'iculo Universitario)}

A medida que los requerimientos de modelamiento se vuelven m\'as complejos, la Tercera Forma Normal (3FN) resulta insuficiente para prevenir ciertas redundancias causadas por claves candidatas superpuestas o dependencias multivaluadas.

\subsection{Forma Normal de Boyce-Codd (BCNF)}
Una relaci\'on $R$ est\'a en BCNF si y solo si, para cada dependencia funcional no trivial $X \rightarrow Y$, $X$ es una \textbf{superclave} de $R$. 

La diferencia cr\'itica con la 3FN radica en que la 3FN permite que $Y$ sea un atributo primo (parte de alguna clave candidata), mientras que BCNF exige estrictamente que el determinante $X$ sea una superclave completa.

\subsection{Cuarta Forma Normal (4FN) y Dependencias Multivaluadas}
Una dependencia multivaluada $X \twoheadrightarrow Y$ indica que para un valor dado de $X$, existe un conjunto bien definido de valores de $Y$ que es completamente independiente de los dem\'as atributos en la relaci\'on.
\begin{itemize}
    \item \textbf{Definici\'on de 4FN:} Una relaci\'on est\'a en 4FN si cumple BCNF y para toda dependencia multivaluada no trivial $X \twoheadrightarrow Y$, $X$ es una superclave.
\end{itemize}

\subsection{Quinta Forma Normal (5FN / PJ-NF) y DKNF}
\begin{itemize}
    \item \textbf{5FN (Project-Join Normal Form):} Trata sobre dependencias de uni\'on (\textit{Join Dependencies}). Una relaci\'on est\'a en 5FN si no puede descomponerse en subrelaciones sin generar tuplas espurias al reconstruirla mediante operaciones \texttt{NATURAL JOIN}.
    \item \textbf{Forma Normal Dominio-Clave (DKNF):} Es el l\'imite te\'orico m\'aximo de la normalizaci\'on. Una relaci\'on est\'a en DKNF si toda restricci\'on de integridad puede deducirse directamente de las definiciones de dominio y de claves.
\end{itemize}

\subsection{Desnormalizaci\'on Estrat\'egica y Rendimiento}
En entornos anal\'iticos y de alta concurrencia de lectura (OLAP), la normalizaci\'on extrema incrementa la sobrecarga de computaci\'on debido al excesivo encadenamiento de \texttt{JOIN}s. La desnormalizaci\'on controlada introduce redundancia deliberada mediante:
\begin{itemize}
    \item Tablas pre-agregadas y tablas sumario.
    \item Columnas calculadas y redundancia de claves for\'aneas.
    \item Vistas Materializadas (\texttt{MATERIALIZED VIEW}) con refresco peri\'odico.
\end{itemize}

---

\section{Casos de Estudio y Pr\'actica Aplicada}

\begin{lstlisting}[style=sqlstyle, caption={Implementaci\'on de Vistas Materializadas para desnormalizaci\'on anal\'itica}]
-- Vista materializada que consolida ventas por cliente y reduce JOINs recurrentes
CREATE MATERIALIZED VIEW mv_customer_sales_summary AS
SELECT
    c.customer_id,
    c.company_name,
    c.country,
    COUNT(DISTINCT o.order_id) AS total_orders,
    SUM(od.unit_price * od.quantity * (1 - od.discount)) AS gross_revenue,
    MAX(o.order_date) AS last_order_date
FROM customers AS c
INNER JOIN orders AS o
    ON c.customer_id = o.customer_id
INNER JOIN order_details AS od
    ON o.order_id = od.order_id
GROUP BY
    c.customer_id,
    c.company_name,
    c.country;

-- Creaci\'on de \'indice \'unico para permitir refrescos concurrentes sin bloquear lecturas
CREATE UNIQUE INDEX idx_mv_customer_sales_pk ON mv_customer_sales_summary (customer_id);

-- Comando de refresco en segundo plano
REFRESH MATERIALIZED VIEW CONCURRENTLY mv_customer_sales_summary;
\end{lstlisting}

---

\section{Taller de Consolidaci\'on del Cap\'itulo}

\begin{businessproblem}
Un sistema de e-learning universitario almacena la asignaci\'on de profesores, cursos y libros de texto recomendados en una sola tabla, violando la 4FN al asociar profesores con cursos y cursos con libros de forma independiente. Se requiere dise\~nar el esquema en 4FN y construir una consulta desnormalizada de alta velocidad para el portal estudiantil.
\end{businessproblem}

\subsection{Soluci\'on T\'ecnica en PostgreSQL}

\begin{lstlisting}[style=sqlstyle, caption={Descomposici\'on en 4FN y consulta consolidada}]
-- 1. Descomposici\'on en 4FN (Eliminaci\'on de dependencias multivaluadas)
CREATE TABLE course_professors (
    course_id VARCHAR(10) NOT NULL,
    professor_id INTEGER NOT NULL,
    PRIMARY KEY (course_id, professor_id)
);

CREATE TABLE course_textbooks (
    course_id VARCHAR(10) NOT NULL,
    isbn VARCHAR(20) NOT NULL,
    textbook_title VARCHAR(150) NOT NULL,
    PRIMARY KEY (course_id, isbn)
);

-- 2. Consulta agregada desnormalizada usando JSONB para el frontend
SELECT
    c.course_id,
    COALESCE(
        jsonb_agg(DISTINCT cp.professor_id) FILTER (WHERE cp.professor_id IS NOT NULL),
        '[]'::jsonb
    ) AS assigned_professors,
    COALESCE(
        jsonb_agg(DISTINCT jsonb_build_object('isbn', ct.isbn, 'title', ct.textbook_title))
        FILTER (WHERE ct.isbn IS NOT NULL),
        '[]'::jsonb
    ) AS required_books
FROM (SELECT DISTINCT course_id FROM course_professors UNION SELECT DISTINCT course_id FROM course_textbooks) AS c
LEFT JOIN course_professors AS cp ON c.course_id = cp.course_id
LEFT JOIN course_textbooks AS ct ON c.course_id = ct.course_id
GROUP BY
    c.course_id;
\end{lstlisting}

\begin{engtip}
\begin{itemize}
    \item \textbf{Agregaci\'on \texttt{JSONB}:} Almacena y proyecta colecciones estructuradas directamente en una sola llamada de base de datos hacia el backend, eliminando problemas de producto cartesiano en clientes HTTP.
\end{itemize}
\end{engtip}
